\documentclass[11pt]{article}
\input{packages.tex}
\usepackage{ulem}
\DeclareMathOperator{\sign}{sign}
\DeclareMathOperator{\Imm}{Im}
\DeclareMathOperator{\Ker}{Ker}
\DeclareMathOperator{\rg}{rg}

% Настройка колонтитулов
\pagestyle{fancy}
\fancyhf{} % Очистить все поля
\fancyhead[L]{Голова — предмет тёмный и исследованию не подлежит}
\fancyhead[R]{\href{https://t.me/ceagest}{Автор}}
\fancyfoot[C]{\thepage}               % Справа внизу

\setlength{\headheight}{13.6pt}

\hypersetup{
    colorlinks=true,
    linkcolor=blue,
    citecolor=blue,
    filecolor=blue,
    urlcolor=blue,
}

\begin{document}

\begin{titlepage}
    \centering
    \vspace*{\fill}  % заполняет пространство сверху

    {\Huge Лекция 10. Евклидовы пространства. Ортогональные базисы и матрицы. Ортогональное дополнение и ортогональные проекции. Метод ортогонализации. Линейные преобразования евклидовых пространств. Преобразование, сопряжённое данному
    \par}

    \vspace*{\fill}  % заполняет пространство снизу
\end{titlepage}

\setcounter{page}{1}

\section{Евклидовы пространства (продолжение)}

\subsection{Ортогональные базисы}

\begin{definition}
    Базис $\{e_1, \ldots, e_n\}$ называется \textbf{ортогональным}, если векторы этого базиса попарно ортогональны, то есть 
    \[(e_i, e_j) = 0 \quad \forall i, j \in \{1, \ldots, n\} : i \neq j\]
    Базис $\{e_1, \ldots, e_n\}$ называется \textbf{ортонормированным}, если он является ортогональным, а также
    \[(e_i, e_i) = 1 \quad \forall i \in \{1, \ldots, n\}\]
\end{definition} \leavevmode
\\
Понятно, что в ортонормированном базисе матрица Грама является единичной, а значит формула для скалярного произведения векторов $x$ и $y$ в ортонормированном базисе имеет вид:
\[(x, y) = x_1 y_1 + \ldots + x_n y_n\]

\begin{theorem}
    Пусть $\{g_1, \ldots, g_n\}$ --- попарно ортогональные ненулевые векторы в $n$-мерном евклидовом пространстве. Тогда они образуют базис, и, кроме того, разложение вектора $x$ по этому базису задаётся формулой
    \begin{equation} \label{eq:1.1}
        x = \sum_{i = 1}^n \frac{(x, g_i)}{\left\lvert g_i \right\rvert^2} g_i
    \end{equation}
\end{theorem}

\begin{proof}
    Действительно, матрица произведений $(g_i, g_j)$ диагональная с ненулевыми элементами на диагонали. Тогда её определитель положителен, а значит по теореме, доказанной в конце предыдущей лекции, имеем линейную независимость векторов набора $\{g_1, \ldots, g_n\}$. При этом количество этих векторов совпадает с размерностью пространства, откуда сразу получаем, что эти векторы составляют базис. \\
    Пусть $x = \alpha_1 g_1 + \ldots + \alpha_n g_n$. Умножая это равенство скалярно на любой из $h_i$, находим, что 
    \[\alpha_i = \frac{(x, g_i)}{\left\lvert g_i \right\rvert^2}\]
    Ясен пень, что отсюда сразу вытекает требуемая формула (\ref{eq:1.1}) разложения вектора $x$.
\end{proof} \leavevmode
\\
Вычислим $(x, x)$ с помощью формулы (\ref{eq:1.1}), работая, естественно, в прежнем ортогональном базисе:
\[(x, x) = \left(\sum_{i = 1}^n \frac{(x, g_i)}{\left\lvert g_i \right\rvert^2} g_i, \sum_{i = 1}^n \frac{(x, g_i)}{\left\lvert g_i \right\rvert^2} g_i\right) = \sum_{i = 1}^n \left(\frac{(x, g_i)}{\left\lvert g_i \right\rvert^2} g_i, \frac{(x, g_i)}{\left\lvert g_i \right\rvert^2} g_i \right) + 2 \sum_{1 \leq i < j \leq n} \left(\frac{(x, g_i)}{\left\lvert g_i \right\rvert^2} g_i, \frac{(x, g_i)}{\left\lvert g_j \right\rvert^2} g_j\right)\]
Поскольку $(g_i, g_j) = 0$ при $i \neq j$, получаем \textbf{равенство Парсеваля}:
\[\left\lvert x \right\rvert^2 = \sum_{i = 1}^n \frac{(x, g_i)^2}{\left\lvert g_i \right\rvert^2}\]

\subsection{Ортогональные матрицы}

Рассмотрим два ортонормированных базиса $e$ и $e' = eS$ в евклидовом пространстве. Если $\Gamma$ и $\Gamma'$ --- матрицы Грама базисов $e$ и $e'$ соответственно, то известна формула
\[\Gamma' = S^T \Gamma S\]
Но базисы являются ортонормированными, а значит $\Gamma' = \Gamma = E$, и формула принимает вид
\begin{equation} \label{eq:par:2:1}
    S^TS = E
\end{equation}
Наоборот, если выполнено условие (\ref{eq:par:2:1}) и исходный базис ортонормированный, то мы получаем $\Gamma' = E$, и новый базис также ортонормированный.

\begin{definition}
    Матрица, удовлетворяющая условию $(\ref{eq:par:2:1})$, называется \textbf{ортогональной} матрицей.
\end{definition} \leavevmode
\\
Как видно из рассуждений выше, ортогональные матрицы и только они могут служить матрицами перехода от одного ортонормированного базиса к другому. Кроме того, равенство (\ref{eq:par:2:1}) равносильно равенству
\[S^T = S^{-1}\]
Из свойств обратной матрицы следует, что
\[SS^T = E\]
Это означает, что матрица $S^T$ также является ортогональной. 

\begin{lemma}
    Произведение $SU$ двух ортогональных матриц $S$ и $U$ --- ортогональная матрица.
\end{lemma}

\begin{proof}
    Действительно,
    \[(SU)^T = U^T S^T = U^{-1} S^{-1} = (SU)^{-1}\]
\end{proof} \leavevmode
\\
Вычисляя детерминант обеих частей равенства (\ref{eq:par:2:1}), мы получим $(\det S)^2 = 1$. Значит, для ортогональной матрицы $\det S = 1$ или $\det S = -1$.

\begin{remark}
    Смекалистый читатель сможет без труда проверить, что любая ортогональная матрица порядка $2$ имеет один из двух видов:
    \begin{equation*}
        \begin{pmatrix}
            \cos \alpha & - \sin \alpha \\
            \sin \alpha & \cos \alpha
        \end{pmatrix}
        \quad 
        \text{ или }
        \quad
        \begin{pmatrix}
            \cos \alpha & \sin \alpha \\
            \sin \alpha & - \cos \alpha
        \end{pmatrix}
    \end{equation*} 
\end{remark}

\subsection{Ортогональное дополнение подпространства}

Пусть $\mathcal{E}_1$ --- $k$-мерное подпространство в $n$-мерном евклидовом пространстве $\mathcal{E}$.

\begin{definition}
    \textbf{Ортогональным дополнением} подпространства $\mathcal{E}_1$ называется множество всех векторов, ортогональных каждому вектору из $\mathcal{E}_1$. Это множество обозначается $\mathcal{E}_1^{\perp}$.
\end{definition}

\begin{theorem}
    Ортогональное дополнение $k$-мерного подпространства в $n$-мерном пространстве есть $(n - k)$-мерное подпространство.
\end{theorem}

\begin{proof}
    Пусть $\{g_1, \ldots, g_k\}$ --- базис в $\mathcal{E}_1$. Вектор $x$ лежит в $\mathcal{E}_1^{\perp}$ тогда и только тогда, когда
    \begin{equation} \label{eq:3.1}
        (x, g_1) = 0, \ldots, (x, g_k) = 0
    \end{equation}
    Действительно, если $x \in \mathcal{E}_1^{\perp}$, то вышеописанные равенства, разумеется, выполнены. Обратно, при выполнении этих равенств $x$ ортогонален любому $y$ из $\mathcal{E}_1$, поскольку
    \[(x, y) = \left(x, \sum_{i = 1}^k y_i g_i \right) = \sum_{i = 1}^k y_i (x, g_i) = 0\]
    Выберем в $\mathcal{E}$ ортонормированный базис и обозначим через $g_{i1}, \ldots, g_{in}$ компоненты вектора $g_i$, где $i \in \{1, \ldots, k\}$, в этом базисе, а через $x_1, \ldots, x_n$ --- компоненты вектора $x$. Условия (\ref{eq:3.1}) запишутся тогда в виде однородной системы из $k$ линейных уравнений с $n$ неизвестными:
    \begin{equation*}
        \begin{cases}
            g_{11} x_1 + \ldots + g_{1n} x_n = 0 \\
            \ldots \ldots \ldots \ldots \ldots \ldots \ldots \ldots \\
            g_{k1} x_1 + \ldots + g_{kn} x_n = 0
        \end{cases}
    \end{equation*}
    Ранг матрицы системы равен $k$, поскольку её строки --- строки из компонент векторов $\{g_1, \ldots, g_k\}$ --- линейно независимы. Таким образом, суммируя все вышедоказанные эквивалентности, заключаем, что множество $\mathcal{E}_1^{\perp}$ определяется однородной системой линейных уравнений ранга $k$, и потому является $(n - k)$-мерным подпространством (в самом деле, вспомните, сколько столбцов имеет фундаментальная матрица однородной системы линейных уравнений ранга $k$). Получили то, что и требовалось, а значит победа постигла нас и теорема доказана!
\end{proof} \leavevmode
\\
Рассмотрим $(\mathcal{E}_1^{\perp})^{\perp}$ --- ортогональное дополнение ортогонального дополнения подпространства $\mathcal{E}_1$. Каждый вектор из $\mathcal{E}_1$ ортогонален каждому вектору из $\mathcal{E}_1^{\perp}$, поэтому $\mathcal{E}_1 \subset (\mathcal{E}_1^{\perp})^{\perp}$. Но $\dim (\mathcal{E}_1^{\perp})^{\perp} = n - (n - k) = k$ по предыдущей теореме. Таким образом, $\mathcal{E}_1$ является подпространством $(\mathcal{E}_1^{\perp})^{\perp}$ и имеет совпадающую с ним размерность, а значит $\mathcal{E}_1 = (\mathcal{E}_1^{\perp})^{\perp}$. \\
Заметим, что $\mathcal{E}_1 \cap \mathcal{E}_1^{\perp} = \{0\}$. В самом деле,
\[\forall x \in \mathcal{E}_1 \cap \mathcal{E}_1^{\perp} \hookrightarrow (x, x) = 0 \iff x = 0\]
Кроме того, по предыдущей теореме сумма размерностей $\mathcal{E}_1$ и $\mathcal{E}_1^{\perp}$ равна $n$. Из этого вытекает 
\begin{lemma}
    Евклидово пространство --- прямая сумма любого своего подпространства и его ортогонального дополнения.
\end{lemma}

\begin{definition}
    Два подпространства $\mathcal{E}_1$ и $\mathcal{E}_2$ называются \textbf{ортогональными}, если $\mathcal{E}_2 \subset \mathcal{E}_1^{\perp}$.
\end{definition} \leavevmode
\\
Заметим, что определение выше в некотором смысле симметрично, так как, если $\mathcal{E}_2 \subset \mathcal{E}_1^{\perp}$, то и $\mathcal{E}_1 \subset \mathcal{E}_2^{\perp}$. Действительно, это простое следствие из очевидного по определению факта:
\[\forall x \in \mathcal{E}_1 \; \; \forall y \in \mathcal{E}_2 \hookrightarrow (x, y) = 0\]

\subsection{Ортогональные проекции}

Так как $\mathcal{E} =  \mathcal{E}_1 \oplus \mathcal{E}_1^{\perp}$, каждый вектор $x \in \mathcal{E}$ однозначно раскладывается в сумму векторов $x_1 \in \mathcal{E}_1$ и $x_2 \in \mathcal{E}_1^{\perp}$. Вектор $x_1$ называется \textbf{ортогональной} проекцией $x$ на $\mathcal{E}_1$. Легко видеть, что $x_2$ --- ортогональная проекция $x$ на $\mathcal{E}_1^{\perp}$. \\
Найдём ортогональную проекцию $x$ на $\mathcal{E}_1$ в предположении, что в $\mathcal{E}_1$ задан некоторый ортогональный базис $\{h_1, \ldots, h_k\}$. Дополним этот базис до ортогонального базиса в пространстве $\mathcal{E}$, присоединив к нему произвольный ортогональный базис $\{h_{k + 1}, \ldots, h_n\}$ из $\mathcal{E}_1^{\perp}$. Так как сумма $\mathcal{E}_1$ и $\mathcal{E}_1^{\perp}$ прямая, искомое разложение вектора $x$ единственно, и мы, группируя слагаемые в формуле (\ref{eq:1.1}), получаем
\begin{equation} \label{eq:4.1}
    x_1 = \sum_{i = 1}^k \frac{(x, h_i)}{\left\lvert h_i \right\rvert^2} h_i
\end{equation}
Если $k = 1$, проекция имеет вид 
\[x_1 = \frac{(x, h)}{\left\lvert h \right\rvert^2} h,\]
и мы видим, что правая часть формулы (\ref{eq:4.1}) --- сумма проекций на ортогональные одномерные подпространства, натянутые на $\{h_1, \ldots, h_k\}$. Так же истолковывается формула (\ref{eq:1.1}), а значит, равенство Парсеваля является обобщением теоремы Пифагора. \\
Наконец, из выполненного по определению условия $(x_1, x_2) = 0$ следует вот такое прекрасное неравенство:
\[\left\lvert x \right\rvert^2 = \left\lvert x_1 + x_2 \right\rvert^2 = \left\lvert x_1 \right\rvert^2 + \left\lvert x_2 \right\rvert^2 + 2(x_1, x_2) = \left\lvert x_1 \right\rvert^2 + \left\lvert x_2 \right\rvert^2 \geq \left\lvert x_1 \right\rvert^2 \]

\subsection{Метод ортогонализации}

Формула (\ref{eq:4.1}) служит основой метода, позволяющего произвольный базис евклидова пространства преобразовать в ортогональный, а затем в ортонормированный. Этот метод называется \textbf{методом ортогонализации Грама-Шмидта}. \\
Пусть в $n$-мерном евклидовом пространстве $\mathcal{E}$ задан некоторый базис $\{h_1, \ldots, h_n\}$. Положим $g_1 = h_1$. Затем из вектора $h_2$ вычтем его ортогональную проекцию на линейную оболочку $g_1$ и положим $g_2$ равным полученной разности:
\[g_2 = h_2 - \frac{(h_2, g_1)}{\left\lvert g_1 \right\rvert^2} g_1\]
Отметим, что $g_2$ раскладывается по $h_1 = g_1$ и $h_2$, причём $g_2 \neq 0$, так как в противном случае $h_1$ и $h_2$ были бы пропорциональны, а значит линейно зависимы. \\
Будем продолжать таким же образом. Допустим, что построены попарно ортогональные ненулевые векторы $\{g_1, \ldots, g_k\}$, причём для любого $i \leq k$ вектор $g_i$ раскладывается по $\{h_1, \ldots, h_i\}$. Положим
\begin{equation} \label{eq:5.1}
    g_{k + 1} = h_{k + 1} - \sum_{i = 1}^k \frac{(h_{k + 1}, g_i)}{\left\lvert g_i \right\rvert^2} g_i
\end{equation}
Вектор $g_{k + 1}$ --- проекция $h_{k + 1}$ на ортогональное дополнение линейной оболочки $\{g_1, \ldots, g_k\}$ (в самом деле, вектор $h_{k + 1}$ единственным образом представим в виде суммы его проекции на линейную оболочку $\{g_1, \ldots, g_k\}$ и его проекции на ортогональное дополнение этой линейной оболочки, что легко следует из доказанного в предыдущей секции равенства $\left\langle g_1, \ldots, g_k \right\rangle \oplus \left\langle g_1, \ldots, g_k \right\rangle^{\perp} = \mathcal{E}$), и потому ортогонален всем $g_i$ при $i < k + 1$. Кроме того, он раскладывается по $\{h_1, \ldots, h_{k + 1}\}$, так как для любого $i \leq k$ вектор $g_i$ раскладывается по $\{h_1, \ldots, h_i\}$. Отсюда следует, что $g_{k + 1} \neq 0$, поскольку иначе векторы $\{h_1, \ldots, h_{k + 1}\}$ оказались бы линейно зависимы. \\
После того как будет преобразован последний вектор $h_n$, мы получим ортогональную систему из $n$ ненулевых векторов. \\
Итак, нами построен ортогональный базис $\{g_1, \ldots, g_n\}$. От него можно перейти к ортонормированному базису $\{e_1, \ldots, e_n\}$ из векторов 
\[e_i = \frac{g_i}{\left\lvert g_i \right\rvert}, \quad i \in \{1, \ldots, n\}\]
Это называется \textbf{нормировкой} базиса $\{g_1, \ldots, g_n\}$. \\
Посмотрим на матрицу перехода $S$ от базиса $\{g_1, \ldots, g_n\}$ к базису $\{h_1, \ldots, h_n\}$. Из равенства $h_1 = g_1$ и формулы (\ref{eq:5.1}) видно, что $h_j$ при любом $j$ раскладывается по $\{g_1, \ldots, g_j\}$, причём его координата по $g_j$ равна $1$. Поэтому элементы $s_{ij}$ матрицы перехода $S$ равны нулю, если они ниже главной диагонали (при $i > j$), и единице при $i = j$ (действительно, это легко следует из того, что столбцы матрицы перехода $S$ есть в сущности координаты векторов нового базиса $\{h_1, \ldots, h_n\}$ в старом базисе $\{g_1, \ldots, g_n\}$). Таким образом, эта матрица --- верхняя треугольная с единицами на главной диагонали.

\section{Линейные преобразования евклидовых пространств}

\subsection{Преобразование, сопряжённое данному}

Всё сказанное ранее о линейных преобразованиях линейных пространств остаётся, конечно, в силе и для евклидовых пространств. С введением скалярного произведения преобразования приобретают новые свойства подобно тому, как векторы приобретают длину.

\begin{definition}
    Линейное преобразование $\varphi^{\ast}$ евклидова пространства называется \textbf{сопряжённым} преобразованию $\varphi$, если для любых векторов $x$ и $y$ этого евклидова пространства имеет место равенство
    \begin{equation} \label{eq:2.1.1}
        \left(\varphi(x), y\right) = \left(x, \varphi^{\ast} (y)\right)
    \end{equation}
\end{definition} \leavevmode
\\
Допустим, что данное преобразование $\varphi$ имеет сопряжённое $\varphi^{\ast}$. Выясним, как связаны матрицы преобразований $\varphi$ и $\varphi^{\ast}$ в некотором базисе $e$. Обозначим эти матрицы через $A$ и $A^{\ast}$. Тогда равенство (\ref{eq:2.1.1}) можно переписать в координатной форме 
\[(Ax)^T \Gamma y = x^T \Gamma A^{\ast} y,\]
где $\Gamma$ --- матрица Грама базиса $e$. Выполнив транспонирование, получаем
\[x^T A^T \Gamma y = x^T \Gamma A^{\ast} y\]
Это равенство показывает, что левая и правая части (\ref{eq:2.1.1}) являются билинейными функциями, а $A^T \Gamma$ и $\Gamma A^{\ast}$ --- матрицы этих функций в базисе $e$. Если значения функций равны при любых $x$ и $y$, то матрицы этих функций равны. Поэтому
\[A^T \Gamma = \Gamma A^{\ast}\]
Итак, матрицы преобразований $\varphi$ и $\varphi^{\ast}$ связаны вышеописанным соотношением. В частности, если базис ортонормированный, то $\Gamma = E$ и
\begin{equation} \label{eq:2.1.2}
    A^{\ast} = A^T
\end{equation}

\begin{theorem}
    Каждое линейное преобразование евклидова пространства имеет единственное сопряжённое преобразование.
\end{theorem}

\begin{proof}
    Для доказательства выберем ортонормированный базис $e$ и рассмотрим линейное преобразование $\psi$, матрица которого в базисе $e$ равна $A^T$. Подставим $\psi$ вместо $\varphi^{\ast}$ в определение (\ref{eq:2.1.1}). Это приведёт к очевидно выполненному матричному равенству $(Ax)^Ty = x^T (A^T y)$. Таким образом, $\psi$ является сопряжённым для $\varphi$, а значит существование показано. \\
    Если бы имелось два преобразования, сопряжённых одному и тому же преобразованию $\varphi$, то в силу (\ref{eq:2.1.2}) их матрицы совпадали бы, а значит единственность также показана и победа вновь постигла нас, товарищи!
\end{proof}

\end{document}